% Responsable: Integrante 1 (tareas 1.1 y 1.2 del plan)
\section{Implementación del lazo cerrado}\label{sec:lazo}
\subsection{Variable controlada y referencia}
Se conserva la planta termoeléctrica de la Tarea 1, con corriente $I$ como entrada y temperatura $T_b$ del bloque frío como salida. Se adopta como objetivo nominal $T_b=5\,{}^\circ\mathrm C$, ambiente de $25\,{}^\circ\mathrm C$ y flujo de aire de $400\,\mathrm{LFM}$. El punto de operación guardado en el proyecto es
\begin{align}
 I_0&=1.190149\,\mathrm A, &T_{b0}&=278.15\,\mathrm K,\\
 T_{h0}&=304.101077\,\mathrm K, &R_{HA}&=1.29\,\mathrm{K/W}.
\end{align}
La referencia representa regulación de una carga térmica, por lo que resulta físicamente coherente con la salida modelada. Para estudiar el seguimiento local se propone un cambio de $5$ a $4\,{}^\circ\mathrm C$ en $t=20\,\mathrm s$. El arranque se estudia por separado, con ambas masas inicialmente a $25\,{}^\circ\mathrm C$ y referencia constante de $5\,{}^\circ\mathrm C$; esta condición está alejada de la linealización.

\subsection{Convención de signo y límites}
Como la ganancia de pequeña señal de la planta es negativa, se define
\begin{align}
 e(t)&=T_{b,\mathrm{med}}(t)-T_{\mathrm{ref}}(t),\label{eq:error}\\
 I(t)&=\operatorname{sat}_{[0,\,2]}\{I_0+K_p e(t)\}.\label{eq:leyP}
\end{align}
Una temperatura superior a la referencia produce un incremento de corriente para $K_p>0$. El sesgo $I_0$ mantiene la corriente necesaria en el equilibrio nominal. El límite superior de $2\,\mathrm A$ se justifica en la Sección~\ref{sec:limite}. La planta física recibe corriente absoluta; la representación lineal recibe $\Delta I=I-I_0$. Las temperaturas internas se expresan en kelvin y se convierten a grados Celsius para la medición y las figuras. La Figura~\ref{fig:lazoP} muestra el lazo resultante.
% NOTA para el integrante 1: los ensayos se prepararon con Imax = 6.4 A.
% Cambiar Imax a 2.0 en parametros_ensayos.mat y repetir (ver Sección límite).

\begin{figure}[!tbp]
\centering
\begin{tikzpicture}[>=Latex,font=\scriptsize,
 b/.style={draw,align=center,minimum height=8mm,inner sep=3pt}]
\node[draw,circle,inner sep=2pt] (e) at (0,0) {$\Sigma$};
\node[b] (k) at (1.35,0) {$K_p$};
\node[b] (s) at (3.0,0) {$+I_0$\\Saturación};
\node[b] (p) at (5.2,0) {Planta\\$I\mapsto T_b$};
\node (y) at (6.5,0) {};
\draw[->] (e)--node[above]{$e$}(k);
\draw[->] (k)--(s);
\draw[->] (s)--node[above]{$I$}(p);
\draw[->] (p)--node[above]{$T_b$}(y.center);
\draw[->] (6.2,0)--(6.2,-1.0)--node[below]{$T_b+n$}(0,-1.0)--node[left]{$+$}(e);
\draw[->] (-0.9,0.65)node[above]{$T_{\rm ref}$}--(-0.9,0)--node[above]{$-$}(e);
\draw[->] (5.2,0.9)node[above]{$Q_{\rm carga}$}--(p);
\end{tikzpicture}
\caption{Lazo de acción inversa con control proporcional. El ruido $n$ se suma en la medición; la perturbación térmica entra en la planta. Fuente: elaboración propia.}
\label{fig:lazoP}
\end{figure}

La implementación propuesta utiliza dos modelos. La red Simscape se reutiliza de la Tarea 1. En la representación lineal se conservan los dos estados térmicos para incorporar condiciones iniciales, una entrada adicional de calor y la temperatura caliente. Una rama con la función de transferencia $G(s)$ verifica la salida de pequeña señal desde equilibrio. Así se mantiene el mismo modelo lineal de la tarea previa y se evita atribuir a una función de transferencia con condiciones iniciales nulas un arranque que no describe por sí sola.
